\odiseachapter{tobal}{Una osada Odisea}\label{cap:osada}

Hace un par de semanas, mi hijo Héctor y yo nos fuimos de tiendas.

Primero fuimos a Zara a comprar unas camisas, es decir, a que él comprara camisas para ambos. Luego hicimos algunos otros recados, preparando su inminente viaje en tren por Europa, el Interrail bendito que une los veranos de tantos jóvenes del continente. Una parte importante de su itinerario pasaba por Grecia, desde donde me mandó ayer unas fotos, sentado con un amigo, frente a la isla de Ítaca. Confieso sin vergüenza que las imágenes me hicieron llorar. Allí estaba mi hijo, tan noble y hermoso como el héroe legendario, allí estaba un sol color escarlata iluminando las costas de un lugar cuyo nombre evoca todo lo que nos es querido, todo lo que nos falta.

Ítaca.

La última tienda que visitamos Héctor y yo fue una librería, ya que tenía el secreto propósito de hacerle el mismo regalo ---la \emph{Ilíada} y la \emph{Odisea}--- que me hizo mi padre a mí hace medio siglo, para que le acompañaran en su viaje a Grecia. Y en efecto, gracias a Nolan, las estanterías rebosaban de ejemplares de los clásicos. Entre las versiones que estuve ojeando, todas ellas primorosamente encuadernadas y con bonitas ilustraciones en las tapas, estaban, cómo no, la venerable traducción en prosa de Segalá y varias traducciones en verso, incluidas las ya clásicas de Fernando Gutiérrez (1951; hoy en Penguin) y José Manuel Pabón (1982, Gredos). Pero la que seleccioné fue la versión de Juan Manuel Rodríguez Tobal, publicada este mismo verano de 2026 en una cuidadísima edición bilingüe de Hiperión.

Y, cómo no, el libro me lo quedé yo (en préstamo). Héctor no necesita leer a Ulises este verano. Es mucho mejor que \emph{sea} Ulises.

¿Pero por qué escogí la traducción más reciente de todas las disponibles? ¿Por qué no conformarme con alguna de las excelentes versiones ya existentes?
El propio Tobal escribe, en un brevísimo prólogo, refiriéndose a su trabajo:

\begin{quotation}
Y aquí he llegado. Y llegar aquí parece que es labor que exige algún argumento.
\end{quotation}

Ciertamente, Tobal, profesor de lenguas clásicas, editor, poeta y reconocidísimo traductor de Safo, Catulo, Ovidio y Anacreonte, entre otros, no se ha caído de ningún guindo. Conoce al dedillo todas las traducciones clásicas al castellano de la \emph{Odisea} y no ignora que atreverse con otra «exige algún argumento». Continúa así:

\begin{quotation}
Claro, la traducción que doy en estas páginas es una, esta sola, una composición particular de entre las infinitas posibilidades que tiene el poema de Homero de ser reproducido, producido de nuevo. Y cientos de personas, de diferentes épocas, edades, lenguas y civilizaciones, se entregaron antes que yo a esta misma labor.
\end{quotation}

Sin duda, hay que ser muy valiente o muy insensato para atreverse.
Pero Tobal no parece especialmente insensato cuando escribe:

\begin{quotation}
¿Debería hacer yo un catálogo de las razones de por qué hice de esta manera lo que hice y no de esa otra? Si me lo impusiera, tendría que argumentar hasta el aburrimiento y la tristeza por qué a lo que tantos dijeron «no se puede» yo digo «sí se puede»: por qué el verso, por qué el hexámetro, por qué la cadencia musical (y no la cadencia de metrónomo), por qué la poesía, por qué la rima, por qué los neologismos, por qué los arcaísmos, por qué la diversidad dialectal, por qué otras lenguas, por qué los nombres parlantes, por qué nunca tan mismos los versos formulares, por qué y dónde y cómo la Ilíada... y muchos más porqués.
\end{quotation}

Nuestro traductor se niega al catálogo y al debate, pero ofrece al perspicaz lector, a la atrevida lectora, claves cruciales para orientarse en el océano que se dispone a abordar. Quizás podamos explorar algunas de ellas juntos.

Empezando por el principio, por el arranque de la \emph{Odisea}. Releamos la traducción en prosa de Segalá:

\begin{quotation}
Háblame, Musa, de aquel varón de multiforme ingenio que, después de
destruir la sacra ciudad de Troya, anduvo peregrinando larguísimo tiempo,
vió las poblaciones y conoció las costumbres de muchos hombres y padeció
en su ánimo gran número de trabajos en su navegación por el ponto, en
cuanto procuraba salvar su vida y la vuelta de sus compañeros á la patria.
Mas ni aun así pudo librarlos, como deseaba, y todos perecieron por sus
propias locuras. ¡Insensatos! Comiéronse las vacas del Sol, hijo de Hiperión; el cual no permitió que les llegara el día del regreso. ¡Oh diosa, hija de
Júpiter!: cuéntanos aunque no sea más que una parte de tales cosas.
\end{quotation}

Y ahora la de Tobal:

\begin{versocita}
De ese hombre, Musa, háblame milerrante, que en tanto extravío\\
dio, tras dejar asolado de Troya el alcázar bendito.\\
Plazas de gentes y gentes él vio, y vio pensares distintos,\\
y por el gran mar tragó en su alma dolores tantísimos\\
su aliento y vuelta amparando de los que llevaba consigo.\\
Pero por más que lo ansiara a salvar no llegó a sus amigos,\\
que por sus propias locuras, las de ellos, cayeron perdidos,\\
¡los insensatos!: las vacas del Sol, de Hiperión Alto el hijo,\\
ellos comieron. Borró día en que uno regresa el dios mismo.\\
Hija de Zeus, de esto ya desde allá donde quieras tú dinos.
\end{versocita}

Antes de seguir, inquieto lector, curiosa lectora, hazme un favor. Lee en voz alta, a ser posible declamando, aún mejor si lo haces en compañía y a orillas del mar. ¿Resuena en tus oídos la cadencia? ¿Sientes el ritmo?

Claro que sí y no necesitas para nada las explicaciones que siguen ---como no las necesitaban los griegos de antaño para escuchar al bardo, sentados en torno a un fuego, quizás con una copa de vino en la mano---. Pero te ruego paciencia, porque vale la pena desvelar los ardides de los que se vale el poeta ---y en este caso me refiero al poeta Tobal traduciendo al Poeta---. Valerosamente, escoge como base de su poema los hexámetros castellanos. Cada hexámetro tiene seis apoyos o ictus. Son las sílabas fuertes que dan el latido, la cadencia del verso. Tomemos como ejemplo la línea que cierra el pasaje que acabamos de leer:

\begin{center}
\small\setlength{\tabcolsep}{3pt}
\begin{tabular}{ccc|ccc|ccc|ccc|ccc|cc}
HI & ja & de & ZEUS & de\_es & to & YA & des & de\_a & LLÁ & don & de & QUIE & ras & tú & DI & nos \\
X & -- & -- & X & -- & -- & X & -- & -- & X & -- & -- & X & -- & -- & X & -- \\
\multicolumn{3}{c|}{D} & \multicolumn{3}{c|}{D} & \multicolumn{3}{c|}{D} & \multicolumn{3}{c|}{D} & \multicolumn{3}{c|}{D} & \multicolumn{2}{c}{E} \\
\end{tabular}
\end{center}

Los seis apoyos del hexámetro ---HI, ZEUS, YA, LLÁ, QUIE y DI--- son las sílabas fuertes que abren cada pie; los grupos de sílabas débiles que siguen a cada uno ---ja de, de\_esto, desde\_a, donde, ras tú y nos--- los completan. Si hay dos débiles, el pie es un dáctilo; si hay una, un troqueo (o un espondeo, si cierra el verso). En este caso, los cinco primeros pies son dáctilos y el último, un espondeo: la composición «ideal», con diecisiete sílabas en total.

¿Son todos los hexámetros de Tobal así de transparentes? El que abre el poema, precisamente, parece saltarse las reglas. Con mis tímpanos de cemento, mi primera reacción fue: ¡anatema! El blasfemo empezaba en una anacrusa (De ese) y la primera tónica no llegaba hasta HOM-bre. Fue el propio autor quien me sacó del error. Como mi amigo y maestro Jenaro Talens\footnote{Corre un rumor, difundido por nadie menos que la moderna reencarnación de Circe, la de las largas trenzas, de que el mayor mérito de la versión inglesa ---la original--- de \emph{Leaving Ithaca} es que dio pie a la formidable traducción de Jenaro, en el volumen bilingüe publicado por Eolas.}, Juanma Tobal tiene oído y por tanto se atreve con dos audaces sinalefas (DÉ-sehom-bre), desplaza el acento de \emph{háblame} hasta la última sílaba (ha-bla-MÉ) y despliega cinco dáctilos impecables rematados por un espondeo:

\begin{center}
\small\setlength{\tabcolsep}{3pt}
\begin{tabular}{ccc|ccc|ccc|ccc|ccc|cc}
DÉ & se\_hom & bre & MÚ & sa\_há & bla & MÉ\,// & mil & e & RRÁN & te & que\_en & TÁN & to\_ex & tra & VÍ & o \\
X & -- & -- & X & -- & -- & X & -- & -- & X & -- & -- & X & -- & -- & X & -- \\
\multicolumn{3}{c|}{D} & \multicolumn{3}{c|}{D} & \multicolumn{3}{c|}{D} & \multicolumn{3}{c|}{D} & \multicolumn{3}{c|}{D} & \multicolumn{2}{c}{E} \\
\end{tabular}
\end{center}

Y aún hay más: tras ese MÉ el verso respira, y esa pequeña pausa corresponde, ni más ni menos, a la cesura pentemímeris del hexámetro griego.\footnote{Del griego \gr{πενθημιμερής}, «de cinco mitades»: si cada pie se divide en dos medios pies, la pentemímeris es la pausa que cae tras el quinto medio pie, es decir, justo después del tercer apoyo. Es el respiro más frecuente del hexámetro homérico, y no lo marca el metro sino la frase: un final de palabra que parte el verso en dos hemistiquios desiguales.}

Todo lo cual queda claro como el agua declamando los versos en voz alta, como te animo a hacer, lírica lectora, bárdico lector.

Por si fuera poco, la fidelidad al original es espectacular. Homero arranca así:

\begin{versocita}
\foreignlanguage{greek}{ἄνδρα μοι ἔννεπε, μοῦσα, πολύτροπον,}
\end{versocita}

Literalmente: \gr{ἄνδρα} (al hombre) --- \gr{μοι} (me) --- \gr{ἔννεπε} (refiere) --- \gr{μοῦσα} (Musa) --- \gr{πολύτροπον}.

A diferencia de Segalá, que arranca con «Háblame, Musa» (asentando un patrón que siguen casi todas las traducciones posteriores de la \emph{Odisea}), Tobal prefiere ceñirse a Homero, que empieza su poema, y no por casualidad, refiriéndose al hombre y no a la musa. De ahí el «De ese hombre, Musa», dejando (como el Poeta) a la deidad detrás del humano.

También nos encontramos en el primer verso el famoso \gr{πολύτροπον} (\emph{polýtropon}), acusativo de \gr{πολύτροπος} (\emph{polýtropos}), una palabra compuesta del prefijo \emph{poly}, muy habitual en Homero, y \emph{tropos}. Una única palabra que nos sirve para ilustrar los desafíos y el encanto de la traducción. \emph{Poly} se traduce como mucho, muy, por encima de..., mientras que \emph{tropos} tiene dos sentidos simultáneos. Por una parte, hacer girar, volver, dar vueltas; por otra, modo, manera, recurso. El traductor al español tiene que elegir. O prima el sentido de «muchos recursos» (como hace Segalá: «varón de multiforme ingenio») o el de vagabundo, como hace Tobal, que introduce el delicioso neologismo \emph{milerrante}.

Lo que nos lleva a otra de las joyas que esta traducción ofrece. Homero abunda, como ya hemos dicho, en \emph{poly}. Tobal, siempre que puede, escoge una traducción directa, prefiriendo, por ejemplo, «milerrante» (donde «mil» traduce un significado común de «poly») a «vagabundo». Y al hacerlo, nos acerca a los griegos de hace tres mil años, no solo porque sus neologismos reproducen fielmente las palabras compuestas del original, sino porque al escucharlos sentimos la \emph{extrañeza} de la poesía, una extrañeza a la que sigue, maravillosamente, el \emph{descubrimiento}. Se diría que cada neologismo que nos encontramos (milerrante, milardido, mientisagaz, milmañas, milpalabrero) es un paisaje extraño a primera vista que se revela familiar de repente, como esos lugares olvidados de la infancia cuya memoria nos viene de nuevo.

Y es que quizás, sin saberlo, estemos recordando el griego de Homero, quizás todos conocemos la \emph{Odisea} de memoria, quizás no somos tan ignorantes como pensamos en estos tiempos iletrados, quizás el acoso vil de las redes sociales y el bombardeo milestúpido de la muchabanalidad no nos han hecho ignorantes sino amnésicos. Y por eso, reconocemos: \gr{πολύτροπος} (polýtropos): milerrante; \gr{περίφρων} (perífron): milpensamientos; \gr{πολύμητις} (polýmētis): milardido, mientisagaz; \gr{πολυμήχανος} (poliméchanos): milmañas; \gr{πολύμυθος} (polýmithos): milpalabrero; \gr{πολύχρυσος} (polýchrysos): milorera; \gr{πολυδένδρεος} (polydéndreos): milarbolado; \gr{πολυμνήστη} (polymnéste): milnovios; \gr{πολυήρατος} (polyératos): bienadamada.

Tobal nos ofrece ese tesoro inagotable que no nos permite leer dos líneas seguidas sin darnos un sobresalto, sin que nos entren ganas de reír y salir a la calle y parar al primer transeúnte que pasa para espetarle a bocajarro: ¡juntanubes, pardediós, bientemplado, onduliebella, machialtivos! ¡Ah, ojiglauca Atenea!

Juan Manuel encaja los mecanismos de composición como un artesano paciente a la vez que un poeta arrebatado. El griego dice de Poseidón que es \gr{ἐνοσίχθων} (enosíchthon o el que hace temblar la tierra). Tobal, según le conviene, usa territremante, territremendo, tanterritremendo, a veces tremante y a veces territodatremante. Su audacia no tiene límite. Su audacia y su oído musical. Me decía en un correo: «me gusta especialmente el efecto terremoto que va en el sonido de la palabra. \gr{γαιήοχος ἐννοσίγαιος} (gaiéochos enosígaios): territodaseñor todaterritremendo». Y yo, a pesar de mi sordera ---herencia, como mi amor por los libros, de mi padre--- creía escuchar temblar la tierra.

Por otra parte, es cierto que la versión de Tobal no es la única que se atreve con hexámetros castellanos. También lo hace José Manuel Pabón, en su clásica traducción, que arranca así:

\begin{versocita}
Musa, dime del hábil varón que en su largo extravío,\\
tras haber arrasado el alcázar sagrado de Troya,\\
conoció las ciudades y el genio de innúmeras gentes.\\
Muchos males pasó por las rutas marinas luchando\\
por sí mismo y su vida y la vuelta al hogar de sus hombres,\\
pero a éstos no pudo salvarlos con todo su empeño,\\
que en las propias locuras hallaron la muerte. ¡Insensatos!\\
Devoraron las vacas del Sol Hiperión e, irritada\\
la deidad, los privó de la luz del regreso. Principio\\
da a contar donde quieras, ¡oh diosa nacida de Zeus!
\end{versocita}

Que estos dos grandes traductores usen hexámetros no es casualidad.\footnote{El hexámetro de Pabón suele describirse, más precisamente, como dactílico anacrústico: cinco acentos en ritmo ternario, con dos sílabas átonas antes del primer acento. Tobal, en cambio, conserva los seis apoyos.} Los dos intentan aproximarse al hexámetro dactílico del original, lo que no es nada fácil. El hexámetro griego es cuantitativo, no acentual. No mide por apoyos (sílabas tónicas), sino por duración, es decir, distingue entre sílabas largas y breves. En otras palabras, en castellano el ictus que marca el ritmo viene dado por el acento. En el griego de la \emph{Odisea}, el ritmo viene marcado por el tiempo que se tarda en pronunciar la sílaba, algo impensable en castellano ---que no distingue vocales largas de breves---, pero también, irónicamente, en griego moderno, que tampoco las distingue. De ahí que la traducción canónica de la \emph{Odisea} al griego moderno, de Kazantzakis y Kakridis, recurra al verso de diecisiete sílabas, muy cercano, de hecho, a los hexámetros de nuestros valerosos traductores.\footnote{Kazantzakis y Kakridis tuvieron que costearse la edición de su \emph{Ilíada}, publicada en 1955. Kazantzakis no llegó a ver impresa su traducción de la \emph{Odisea}, que apareció póstumamente en 1965.}

Entonces, ¿en qué se diferencian Pabón y Tobal?
El texto de Pabón es espectacular y una referencia obligatoria, incluso para el propio Tobal. Pero Juanma\footnote{Habrán notado la observafiel lectora, el veolotodo lector, que me permito la libertad de llamar a Juan Manuel por su nombre, echando mano de la cercanía de la amistad y reservando el «don» para sus clásicos ---y difuntos--- predecesores.} no se ha tirado unos cuantos años traduciendo la \emph{Odisea} para no tener nada que decir. Leamos el siguiente pasaje del impecable don José Manuel.

\begin{versocita}
Cuantos antes habían esquivado la abrupta ruina,\\
en sus casas estaban a salvo del mar y la guerra;\\
solo a él, que añoraba en dolor su mujer y sus lares,\\
reteníale la augusta Calipso, divina entre diosas,\\
en sus cóncavas grutas, ansiosa de hacerlo su esposo.
\end{versocita}

Y aquí la de Tobal, buscareyertas:

\begin{versocita}
Ya los demás, los que habían esquivado la muerte escarpada,\\
eran en casa por fin, bien a salvo de mar y batalla.\\
Huérfano así de regreso y de esposa, solo a él lo guardaba\\
ninfa divina entre diosas, Calipso la soberana\\
dentro de cuéncavas grutas por ver si con él se casaba.
\end{versocita}

Observa, ojiglauca lectora, divinal lector, la diferente naturaleza de estas dos hermosas bestias, en apariencia tan similares. La primera se puede leer mentalmente o en voz alta y no hay mucha diferencia. Los versos entran bien y el ritmo es impecable, quizás demasiado impecable. Hay algo de predecible en ellos. Su escritura pulcra, sus cláusulas exactas, su métrica invariable tienden a causar el efecto de que, al leer en voz alta, cada verso caiga con el mismo patrón, la misma duración y el mismo golpe final. Leído en silencio, el texto de Pabón fluye como un río. Al declamarlo, se vuelve un metrónomo.

Juanma, sacándose sinalefas de la manga como un tahúr, pide a gritos que le declamemos. En cada verso, a pesar de su rigor métrico y su fidelidad al griego, esconde alguna carta marcada. Donde Pabón dice «abrupta ruina», él nos ofrece una violenta «muerte escarpada». Frente al castellano impecable «en sus casas estaban a salvo del mar y la guerra» escribe «eran en casa por fin, bien a salvo de mar y batalla», una construcción arcaizante, pero cercana al original. Frente al educadísimo «solo a él, que añoraba en dolor su mujer y sus lares», replica con «Huérfano así de regreso y de esposa, solo a él lo guardaba». ¡Huérfano de regreso y esposa! Imposible no dar un brinco al leerlo, imposible no recordar al muchacho de entonces. Medio siglo después de leer la \emph{Odisea} por primera vez con mi primo Pedro, vuelvo a sentirme poeta (y niño) de nuevo.
